\cleardoublepage

\section{Introduction}

% Im Einsatz unter Chemikalienschutz ist jede Sekunde
% entscheidend – die Träger:innen müssen wichtige
% Informationen wie Atemvorrat, Innendruck, oder Temperatur
% sofort erfassen können.
% innovative Konzepte zur Darstellung der Daten im Sichtfeld
% visuelle Klarheit, schnelle Wahrnehmbarkeit und ergonomische Gestaltung
% • Sammlung relevanter Informationen beim Einsatz von
% Ganzkörperschutzbekleidung sowie deren Gewichtung
% • Konzepte der Datenvisualisierung unter Beachtung der
% Wahrnehmbarkeit und der Nachvollziehbarkeit
% • Eruierung möglicher Technologien zur Visualisierung

Despite rapid advances and numerous potential applications for augmented reality, it has not yet been widely adopted in industrial settings.
One such area of application is in activities where chemical protective clothing is needed. This inherently dangerous work involves significant physical exertion.
These activities require constant awareness of the environment and access to information, which is provided by monitoring devices and radio communication in current practice.

The protective equipment restricts the senses and limits mobility, which makes every movement a challenge.
A head-up display is able to integrate relevant information into the field of view, thereby enabling hands-free and immediate access, potentially replacing current methods.

However, workers are already pushed to their limits by the equipment and task demands. The HUD must not introduce additional burden—neither in terms of weight nor cognitive load—and it must be possible to utilize the system without interrupting current tasks.



First, the requirements of the intended use case, the limitations imposed by AR technology, and biological constraints are identified. 
Next, the mechanisms of attention are introduced, including how they can be utilized to design a head-up display capable of providing relevant information without causing distractions.
This is followed by a compilation of design options that enable rapid access to and comprehension of information within the aforementioned constraints.
Finally, design guidelines are proposed that compile the results of studies on AR and attention management and are illustrated using a potential HUD design.


% An established approach for monitoring information without distraction is the implementation of peripheral displays.
% As the name implies, they present information in the peripheral vision, enabling passive monitoring.
